\documentclass[10pt]{article}
\usepackage[margin=2.5cm]{geometry}
\usepackage{booktabs}
\usepackage{float}
\usepackage{graphicx}

\title{Labwork 3\\\large HPC -- RGB to grayscale}
\author{
Le Nghiem Thanh Thuy - 2540023\\
University of Science and Technology of Hanoi
}
\date{}

\begin{document}
\maketitle

\section{Preparation}
The image is \texttt{astronaut()} from \texttt{skimage.data}, size $512\times512$ with 3
channels, so it has 262144 pixels. Before computing anything I reshape the image from
$(512, 512, 3)$ to $(262144, 3)$, so the image becomes a flat list of pixels and one
thread can take care of exactly one pixel.

The grayscale value of a pixel is the average of its three channels,
$g = (R + G + B) / 3$, and the same value is written back to all three output channels.

\section{Results}
The two versions give exactly the same image (\texttt{np.array\_equal} returns
\texttt{True}).

\begin{table}[H]
\centering
\caption{Execution time for the grayscale conversion of a $512\times512$ image
(Tesla T4 on Google Colab).}
\label{tab:time}
\begin{tabular}{lrr}
\toprule
\textbf{Version} & \textbf{Time (ms)} & \textbf{Speedup} \\
\midrule
CPU & 257.92 & 1.0x \\
GPU & 3.92 & 65.8x \\
\bottomrule
\end{tabular}
\end{table}

\begin{figure}[H]
\centering
\includegraphics[width=0.45\textwidth]{original.png}
\includegraphics[width=0.45\textwidth]{grayscale.png}
\caption{The original image (left) and the grayscale result (right).}
\label{fig:result}
\end{figure}

\section{Different block sizes}
I ran the same kernel with block sizes 16, 32, 64, 128, 256, 512 and 1024.
The block size does not change the image, only how the threads are grouped, so the
result is the same every time. Here I only measured the kernel and kept the image on
the device, because the memory copies are much slower than the kernel and would hide
the difference. Each block size is launched once as a warm-up and then 100 times, and
I take the average.

\begin{figure}[H]
\centering
\includegraphics[width=0.8\textwidth]{blocksize.png}
\caption{Kernel time for different block sizes.}
\label{fig:blocksize}
\end{figure}

A block size of 16 is clearly the slowest (0.152 ms). A warp is 32 threads, so with only
16 threads per block half of every warp does nothing. From 32 upwards all the times are
between 0.084 and 0.104 ms, which is quite close to each other, and the fastest one here
is 64. The kernel is very simple and each thread only reads 3 bytes and writes 3 bytes,
so the block size does not matter much as long as it is a multiple of the warp size.

\section{Comment}
The GPU is much faster than the CPU loop.
The first GPU run is slower because it also has to compile the kernel and set up CUDA
(148.44 ms, against 3.92 ms for the next one), so I only measured the time after a
warm-up run.

\end{document}
